\documentclass[11pt,a4paper]{article}
\usepackage[margin=2.5cm]{geometry}
\usepackage{amsmath,amssymb}
\usepackage{booktabs}
\usepackage{natbib}
\usepackage{graphicx}
\usepackage[colorlinks=true,linkcolor=blue,citecolor=blue,urlcolor=blue]{hyperref}

\newcommand{\nh}{N_{\mathrm{H}}}
\newcommand{\kev}{\,\mathrm{keV}}
\newcommand{\kms}{\,\mathrm{km\,s^{-1}}}
\newcommand{\zsun}{Z_\odot}
\newcommand{\code}[1]{\texttt{#1}}

\title{Inference and evaluation with the SpexAI emulator:\\
forward model, absorption, temperature distributions, and bias testing\\
\large Methodology report (companion to the emulator technical report)}
\author{D.~Huppenkothen \and Claude (Anthropic)}
\date{2026 July}

\begin{document}
\maketitle

\begin{abstract}
This report documents the inference side of \code{spexai}: how a full observed
X-ray spectrum is assembled from the per-element neural operators, how Galactic
absorption and velocity broadening are applied, how we fit single- and
multi-temperature (differential emission measure, DEM) models with MCMC and
nested sampling, and how we test the emulator for bias against an independent
SPEX-derived ground truth. It also reports a targeted study of whether the
emulator's residual emissivity error on the hot iron-peak metals---which miss
the $0.1\%$ training target---biases the quantities XRISM actually measures for a
Perseus-like cluster (Sec.~\ref{sec:hotfloor}). It is deliberately kept separate
from the emulator \emph{development} report, which covers architecture and
training. Design
choices flagged here as approximations (PCHIP ground truth, native-resolution
truth broadening, \code{wabs} vs.\ \code{tbabs} cross-sections) are called out
explicitly so they can be revisited.
\end{abstract}

\section{Forward model}
An observed spectrum is built as
\begin{equation}
C_c = t_{\rm exp}\,N \sum_e R_{ec}\, A_e
      \left[\sum_Z a_Z\, f_{Z}\!\big((1{+}z)E_e; T, v\big)\right]
      \, e^{-\nh\,\sigma(E_e)} ,
\end{equation}
where $f_Z$ is the emulated per-element flux, $a_Z$ the solar-relative
abundance, $R$/$A$ the RMF/ARF, $z$ the source redshift, $v$ the velocity
broadening, and $e^{-\nh\sigma}$ the Galactic absorption screen. The joint
emitter is \code{JointOperatorModel}; folding is \code{Response.fold}. Because
folding is linear, a temperature distribution is summed \emph{before} folding
(Sec.~\ref{sec:dem}).

\section{Velocity broadening}
Broadening is a Doppler redistribution, a photon from $E_j$ landing at
$E_j(1+v_{\rm los}/c)$ with $v_{\rm los}\sim\mathcal N(0,v^2)$. We use a
\emph{hybrid} scheme: the smooth continuum is convolved by FFT in uniform
log-energy (shift-invariant to first order in $v/c$, differentiable, batched
over per-spectrum velocities), while line-head lines are deposited as exact
bin-integrated Gaussians (\code{erf} differences) at their exact energies. The
exact redistribution matrix \code{direct\_broadening\_matrix} is retained as the
reference; the deployed FFT+analytic path agrees with it to well below the
emulator's own error.

\section{Why an operator, given that broadening needs a fine grid}
It is reasonable to ask why we train a grid-portable operator if the broadened
forward pass still routes the continuum through a fine log-uniform grid. The
fine grid is a property of the \emph{FFT broadening method} (its resolution is
set by the kernel width $\sigma/E=v/c$), not of the emulator: a fixed-grid model
would need it too, and worse---upsampling already-binned fluxes smears lines.
The operator's payoff in the broadened pipeline is therefore twofold: (i)
\emph{analytic, grid-exact lines}, deposited directly onto the instrument grid
and never passed through a lossy rebin, which is exactly the
resolution-critical part; and (ii) \emph{smooth, batched temperature
dependence}, which is what makes DEM integration over hundreds of temperatures
cheap and differentiable---a per-SPEX-temperature table would require PCHIP
interpolation at every likelihood call. Grid portability is fully realised for
the line head and for the unbroadened path (\code{forward\_on\_grid}); the exact
matrix route even broadens directly on the instrument grid with no fine grid,
and is kept as the reference.

\section{Galactic absorption}
Photoelectric absorption multiplies the source flux by
$T(E)=e^{-\nh\sigma(E)}$ in the \emph{observed} frame. The cross-section per
hydrogen atom $\sigma(E)$ is taken from an established model, not invented:

\begin{itemize}
\item \textbf{\code{tbabs}} \citep*{wilms2000} is tabulated once from XSPEC
(PyXspec, with \code{abund wilm}), inverted to $\sigma(E)$ and cached as a tensor
(\code{scripts/build\_tbabs\_table.py}); runtime never calls XSPEC. conda-forge
ships no XSPEC for macOS, but HEASoft is installable via the HEASARC conda
channel, from which the table was generated
($\sigma(1\kev)=1.74\times10^{-22}\,\mathrm{cm^2}$). \code{Absorption.default()}
uses it when the cached table is present.
\item \textbf{\code{wabs}} \citep*{morrison1983} is the dependency-free fallback,
an analytic polynomial validated against its published coefficients
($\sigma(1\kev)=2.422\times10^{-22}\,\mathrm{cm^2}$), used when the \code{tbabs}
table is absent (e.g.\ a fresh checkout without HEASoft).
\end{itemize}

\paragraph{Insertion point (instrument-resolution independence).}
Absorption acts on flux \emph{density}, so applying $T(E_{\rm centre})$ to an
already-rebinned coarse spectrum would mis-integrate across a bin and smear the
photoelectric edges, differently per instrument. Instead $T$ is applied on the
native fine grid \emph{before} rebinning (continuum) and at the exact line
energies (lines), evaluated at the observed energy $E_{\rm rest}/(1{+}z)$. The
coarse RMF then only ever sees an already-absorbed, flux-conserved spectrum, so
the result is instrument-resolution independent by construction and needs no
operator retraining. The same screen multiplies the emulator and the SPEX
truth, so the bias round-trip stays self-consistent regardless of the
cross-section choice.

\section{Absolute normalisation and distance}
The ``norm'' fed to \code{predict\_counts} is the SPEX \code{cie} normalisation:
the emission measure $Y \equiv n_H n_e V$ in units of $10^{64}\,\mathrm{m^{-3}}$.
The per-element operators were trained on SPEX spectra at the SPEX reference
$Y=1$ and reference luminosity distance $D_{\rm ref}=10^{22}\,\mathrm{m}$ (the
SPEX default \code{distance} unit), so physical detector counts are
\begin{equation}
C_c = t_{\rm exp}\; Y \left(\frac{D_{\rm ref}}{D}\right)^{\!2}
      \; u_{\rm A}\; \sum_e R_{ec} A_e\, f_e ,
\end{equation}
with $u_{\rm A}=10^{-4}$ converting SPEX's SI photon flux
($\mathrm{ph\,s^{-1}\,m^{-2}\,bin^{-1}}$) to the per-$\mathrm{cm^2}$ convention of
the OGIP ARF. These constants live in \code{spexai.inference.units}. $Y$ and $D$
are perfectly degenerate in a single spectrum, so $D$ is supplied fixed (from the
known redshift/distance) and $Y$ is fitted.

\paragraph{Validation against SPEX.} We installed SPEX (SPEXACT 2.07, native
Apple-Silicon conda build) and compared the emulator's \emph{absolute} flux to a
SPEX \code{cie} spectrum at $T=4\kev$, $Y=1$, $D=D_{\rm ref}$
(\code{scripts/validate\_spex\_norm.py}). The emulator reproduces SPEX to a
median $\sim2.5\%$ across the $3$--$8\kev$ continuum (ratio $0.99$ at $2$--$5\kev$,
falling to $0.95$ at $8\kev$). This confirms the flux units
($\mathrm{ph\,s^{-1}\,m^{-2}\,bin^{-1}}$, hence $u_{\rm A}=10^{-4}$), $Y_{\rm
ref}=1$ and $D_{\rm ref}=10^{22}\,\mathrm{m}$ -- there is no order-of-magnitude
unit error. The mild high-energy deficit is the expected signature of the
present model covering only 16 of 30 elements (SPEX's additional metals add
continuum), and should shrink with the full 30-element model.

\section{Independent SPEX ground truth}\label{sec:truth}
To test the emulator for bias we need spectra \emph{not} produced by it.
\code{SpexTruthModel} interpolates the preprocessed per-element SPEX training
spectra with per-bin monotone cubic (PCHIP) in $\log_{10}T$, sums elements
linearly by abundance, broadens with the exact redistribution matrix on the
native SPEX grid, and folds through the response. This is a documented
approximation to a live SPEX call: the cache carries no abundance dimension
(abundances enter only through the linear CIE sum, which is exact), broadening
is at the native SPEX resolution, and off-grid temperatures come from PCHIP
rather than SPEX. Against the emulator it agrees to $\sim0.1$--$0.7\%$ in flux
for iron, i.e.\ both approximate SPEX at the sub-percent level; the residual is
what the bias study probes.

\section{Temperature distributions and abundances}\label{sec:dem}
A DEM is a weighted sum $\sum_g w_g\, f(T_g)$ over a temperature grid, evaluated
by \code{predict\_counts\_dem} (folding once, after the sum). Shapes are drawn
from standard libraries: Gaussian in $\log_{10}T$ or in $T$, a two-component
mixture, and a log-normal, all thin wrappers over \code{scipy.stats}; the
non-parametric \code{BinnedDEM} carries one free, normalised weight per coarse
bin with no smoothness prior. Abundances use \code{AbundanceModel}: a global
metallicity scaling all metals, with individual elements freed or tied to a
fraction of iron (an $[X/\mathrm{Fe}]$ ratio), reproducing the thesis scheme
(global $Z$ with free N/Fe, O/Fe) in a reusable object.

\section{Bias testing}
Two distinct questions are easily conflated here, and we separate them
deliberately. \emph{Is the posterior sampler correct for the model being
fitted?} is answered by simulation-based calibration (SBC): parameters are drawn
from the prior, data are simulated \emph{from the fitted model itself}, and the
rank of the truth among the posterior draws is recorded; a uniform rank
histogram indicates calibration. \emph{Is the emulator unbiased with respect to
SPEX?} is a different question, and SBC is the wrong instrument for it: a rank
statistic carries no physical units, has no exposure axis, and cannot separate
emulator error from sampler error, since all sources of misspecification produce
the same non-uniformity. It is also the most expensive available probe, costing
a full posterior per simulation.

We therefore use SBC as the \emph{control}---it establishes that the sampler is
calibrated, so that any residual offset in the studies below is attributable to
the emulator---and measure emulator bias directly with the Fisher estimator of
Sec.~\ref{sec:fisherbias}, which reports the offset in physical units and in
units of the statistical error at a stated exposure. Ranks computed on SPEX-
injected data remain useful as a cheap global screen over the whole prior
volume, since they aggregate over regions a grid might miss, but they are a
misspecification test rather than a calibration test and are labelled as such.

Where posterior fits are used for bias---to confirm points the Fisher screen
flags, and to validate the linearisation---observations are simulated from
\code{SpexTruthModel}, fitted with the emulator (\code{run\_emcee};
\code{run\_ultranest} for evidence/spot checks), and scored by the
median$-$truth pull and central-interval coverage. An emulator-self round-trip
is the control that separates emulator bias from sampler and statistical
scatter. One caveat applies to all of these: \code{SpexTruthModel} interpolates
the same cached SPEX rows the emulator was trained on, so on-grid it reproduces
SPEX exactly while off-grid it is independent of the emulator's \emph{function}
but not of its \emph{training data}. With a PCHIP interpolation error of
${\sim}0.003\%$ against an emulator error of ${\sim}0.3\%$ this is a safe
margin, but it measures interpolation fidelity rather than extrapolation.
Parameter ranges (Table~\ref{tab:ranges}) are
literature-grounded; note that real clusters sit near $Z\sim0.3\,\zsun$, well
below the thesis's fiducial $\sim1\,\zsun$, so the calibration must span that
regime. A \emph{Perseus} showcase (single-temperature core and a Gaussian DEM)
uses published values ($z=0.0179$, $\nh\sim1.4\times10^{21}\,\mathrm{cm^{-2}}$,
$kT\sim3.9\kev$, $Z_{\rm Fe}\sim0.55\,\zsun$, $\sigma_v\sim180\kms$) folded
through XRISM/Resolve, as a concrete recovery demonstration.

\begin{table}[t]
\centering
\caption{Parameter ranges. ``Realistic'' is the primary SBC volume; ``Extreme''
is the stress-test extension. $kT$ is capped at the $0.2$--$10\kev$ training
band.}
\label{tab:ranges}
\begin{tabular}{lccc}
\toprule
Parameter & Realistic & Extreme & Perseus \\
\midrule
$kT$ (keV)        & 1--8       & 0.3--10   & $\sim$3.5--4 \\
$Z$ ($\zsun$)     & 0.2--0.7   & 0.1--1.5  & $\sim$0.5--0.6 \\
$[\alpha/\mathrm{Fe}]$ & 1.0--1.5 & 0.5--2 & $\sim$1 \\
$v$ (km/s)        & 0--300     & 0--600    & $\sim$150--200 \\
$\nh$ (cm$^{-2}$) & $5{\times}10^{19}$--$3{\times}10^{21}$ & $\to10^{22}$ & $\sim$1.4$\times10^{21}$ \\
$z$               & 0.01--0.3  & $\to$1    & 0.0179 \\
\bottomrule
\end{tabular}
\end{table}

\section{Systematic bias without fitting: the Fisher estimator}\label{sec:fisherbias}
This section derives the estimator used both for the hot-element study
(Sec.~\ref{sec:hotfloor}) and for the parameter-space sweep
(Sec.~\ref{sec:biassweep}). It answers a specific question: \emph{if a real
spectrum is fitted with the emulator rather than with SPEX, how wrong are the
recovered parameters, and at what exposure does that matter?} The estimator
obtains both the systematic offset and the statistical error from two
\emph{noise-free} spectra, with no fitting and no simulated observations. Since
that is counter-intuitive, we set it out in full.

\subsection{Definitions}
Throughout, $c$ indexes instrument channels inside the fit band, $i,j$ index the
$n$ free parameters, and $\theta_0$ is the true parameter vector at which
everything is evaluated.

\begin{description}
\item[$\theta$] the free parameter vector, $n$ components. For the literature
strategy these are the eight free abundances (Si, S, Ar, Ca, Cr, Mn, Fe, Ni),
the temperature descriptor ($kT$, or $T_{\rm mean}$ and $T_{\rm sigma}$ for a
Gaussian DEM), the velocity dispersion $\sigma_v$, the absorbing column $\nh$,
and $\log_{10}$ of the normalisation, so $n=12$ (13 for the DEM).
\item[$d_c$] the \emph{noise-free} expected counts in channel $c$ predicted by
the independent SPEX ground truth (\code{SpexTruthModel}, Sec.~\ref{sec:truth}),
evaluated at $\theta_0$ and folded through the instrument response. This is the
``correct answer'': what nature would produce, in the absence of counting
noise. It involves no emulator.
\item[$\mu_c \equiv \mu_c(\theta_0)$] the noise-free expected counts predicted
by the \emph{emulator} at the same $\theta_0$, through the same response. If the
emulator were exact, $\mu_c=d_c$ for every channel.
\item[$J_{ic} \equiv \partial\mu_c/\partial\theta_i$] the Jacobian of the
emulator: how much channel $c$ responds to a unit change in parameter $i$.
Computed by central finite differences, $2n{+}1$ emulator evaluations in total.
\item[$r_c \equiv d_c/\mu_c - 1$] the fractional emulator error in channel $c$.
This is a \emph{deterministic} quantity---a property of the emulator, not of any
observation---and it is the only place the truth enters.
\item[$N_{\rm ref}$] a reference in-band count level ($10^5$ in the code) at
which $\sigma_{\rm stat}$ is quoted; results are rescaled analytically to any
other exposure via Eq.~\eqref{eq:scaling}.
\end{description}

\subsection{Two noise-free spectra, two different jobs}
The calculation proceeds in two stages, which in the code
(\code{scripts/bias\_sweep.py}) are separate because $d$ does not depend on the
emulator and is therefore reusable across retrained models.

\emph{Stage 1 (truth)} computes $d$: pure SPEX physics at $\theta_0$, summed
over elements, broadened, absorbed and folded. \emph{Stage 2 (bias)} computes
$\mu$ and $J$ from the emulator, and combines them with $d$ as below.

It is worth being explicit about why neither stage needs noise, because the two
spectra play different roles: $d$ supplies the \emph{systematic} (through $r$),
while $\mu$ supplies the \emph{variance} (analytically, as shown next).

\subsection{Where the statistical error comes from}
For Poisson data the variance is not an independent quantity---it is fixed by
the mean:
\begin{equation}
\mathrm{Var}(d_c) = \mathbb{E}[d_c] = \mu_c .
\end{equation}
A channel expecting $10^4$ counts has $\sigma=100$; one expecting $10^2$ has
$\sigma=10$. A noise-free model spectrum therefore already encodes the complete
noise structure of an observation that was never simulated. No random draw is
required, and drawing one would only add scatter to be averaged away again.

Concretely, the Poisson log-likelihood (dropping the $\theta$-independent
$\log d_c!$) is
\begin{equation}
\log L(\theta) = \sum_c \big[\, d_c \log\mu_c(\theta) - \mu_c(\theta) \,\big],
\end{equation}
and the Fisher information is its expected curvature at $\theta_0$. Taking the
expectation with $\mathbb{E}[d_c]=\mu_c$ makes the second-derivative term vanish
and leaves
\begin{equation}
F_{ij} \;=\; \mathbb{E}\!\left[-\frac{\partial^2 \log L}
{\partial\theta_i\,\partial\theta_j}\right]
\;=\; \sum_c \frac{J_{ic}\,J_{jc}}{\mu_c}
\;=\; \sum_c \frac{J_{ic}\,J_{jc}}{\mathrm{Var}(d_c)} .
\label{eq:fisher}
\end{equation}
The final form is the interpretive one: $F$ is an
\emph{inverse-variance-weighted inner product of the derivatives}, exactly the
weighting a $\chi^2$ fit would use, with the Poisson variances supplied by the
model itself. $J_{ic}$ says how strongly channel $c$ responds to parameter $i$;
$1/\mu_c$ says how much that channel can be trusted. A parameter that moves many
bright channels is well constrained; one that moves a few faint channels is not.
In code this is the single line \code{F = (J / mu0) @ J.T}.

The Cram\'er--Rao bound then gives the covariance of any unbiased estimator,
\begin{equation}
\mathrm{Cov}(\hat\theta) \;\ge\; F^{-1},
\qquad
\sigma_{{\rm stat},i} = \big[(F^{-1})_{ii}\big]^{1/2}.
\end{equation}
The conceptual point is that $F^{-1}$ is not the error of one particular noisy
observation: it is the error \emph{averaged over all noise realisations}, an
expectation rather than a measurement. That is precisely the quantity wanted
here---whether a bias would be detectable in a typical observation, not whether
it happened to be detectable in one draw---and it is more stable than estimating
the same thing from a finite number of simulated fits.

\subsection{Where the systematic comes from}
Because the emulator error is small ($|r_c|\ll1$; see the residual
characterisation in Sec.~\ref{sec:hotfloor}), the maximum-likelihood solution
obtained by fitting $d$ with the emulator lies one Newton step from $\theta_0$.
The score (gradient of $\log L$ at $\theta_0$, with data $d$) is
\begin{equation}
s_i = \frac{\partial \log L}{\partial\theta_i}\bigg|_{\theta_0}
    = \sum_c \Big(\frac{d_c}{\mu_c}-1\Big) J_{ic}
    = \sum_c r_c\, J_{ic},
\label{eq:score}
\end{equation}
which projects the channel-space emulator error onto each parameter direction,
and the Newton step is
\begin{equation}
b_{\rm sys} = F^{-1} s .
\label{eq:bsys}
\end{equation}
$b_{\rm sys}$ is the \emph{parameter-space bias}: where the fit would converge
minus where it should, in physical units (dex of abundance, keV of temperature,
$\kms$ of velocity). This is the step that converts a spectrum-space accuracy
statement---the mean relative error quoted per element in the development
report---into the quantity the science actually publishes. A perfect emulator
gives $r=0$, hence $s=0$ and $b_{\rm sys}=0$, which is the natural sanity check.

Note that Eq.~\eqref{eq:fisher} evaluates the variance weight with the
emulator's $\mu_c$ rather than the truth $d_c$. The two differ by $|r_c|$, a
fraction of a percent, which is negligible inside a variance weight; this is the
$d\simeq\mu$ approximation noted in the code.

\subsection{The crossover, and why one exists}
Suppose the exposure (or source flux) is scaled by a factor $k$, so that
$\mu_c \to k\mu_c$ and hence $J_{ic} \to kJ_{ic}$. Then from
Eqs.~\eqref{eq:fisher}--\eqref{eq:bsys}:
\begin{align}
F &\to kF, & \sigma_{\rm stat} &\to k^{-1/2}\sigma_{\rm stat}
   &&\propto N^{-1/2}, \nonumber\\
r &\to r, & s &\to ks, & b_{\rm sys} = F^{-1}s &\to b_{\rm sys}
   &&\text{(unchanged).}
\label{eq:scaling}
\end{align}
The fractional residual $r_c=d_c/\mu_c-1$ is a \emph{ratio}, so it is invariant
under the rescaling; the factors of $k$ in $s$ and $F^{-1}$ then cancel exactly.

This asymmetry is the entire basis of the test. The emulator's bias is a
\emph{fixed offset in physical units} that no amount of observing time reduces,
while the statistical error shrinks as $N^{-1/2}$. They cross at
\begin{equation}
N_\star = N_{\rm ref}\,\big(\sigma_{\rm stat}/|b_{\rm sys}|\big)^2 ,
\end{equation}
the in-band count level at which the two are equal. Below $N_\star$ the emulator
error is invisible beneath the counting noise; above it, the analysis reports a
systematically shifted value with confident-looking error bars. Reporting
$N_\star$ rather than a goodness score is what makes the result actionable: it
is directly comparable to the counts in a real observation.

For shape parameters $b_{\rm sys}$ is additionally independent of the absolute
count scale because the normalisation parameter absorbs it, so $N_\star$ is a
fixed property of the triple (emulator, modelling strategy, temperature model).

\subsection{Cost and domain of validity}
Both $\sigma_{\rm stat}$ and $b_{\rm sys}$ are read off the \emph{same}
finite-difference Jacobian, so the total cost is $2n{+}1$ emulator evaluations
(25 for $n=12$) plus one truth spectrum---no posterior sampling anywhere. This
is what makes a sweep over hundreds of parameter combinations affordable when a
single MCMC posterior costs hours (Sec.~\ref{sec:biassweep}).

The price is that Eq.~\eqref{eq:bsys} is first order in $r$. The linearisation
was cross-checked against converged GPU MCMC fits at several count levels
(Sec.~\ref{sec:hotfloor}): the recovered offsets agree in direction (10 of 11
parameter signs) but can reach $1.5$--$3\times$ the linear prediction at
$10^8$ counts, where the fit is far into the systematics-dominated regime.
The estimator should therefore be read as a calibrated \emph{screen and
ranking}---reliable for deciding which parameter combinations deserve
attention, and conservative by a factor of a few in magnitude---rather than as a
final number. Points it flags are then confirmed with full posteriors.

\section{Practical impact of the hot-element emissivity floor}\label{sec:hotfloor}
The per-element operators are trained toward a relative emissivity error of
$0.1\%$, a threshold set for high signal-to-noise XRISM spectra rather than
derived from a science requirement. Seven hot, line-rich iron-peak metals miss
it (development report): Cr, Mn and Ni are marginal-to-failing and Ti, Co, Cu and
Zn fail outright, on a line-head temperature-conditioning floor whose signature
is a near-uniform $\sim0.3$--$0.4\%$ \emph{multiplicative} error on line-bin
amplitudes. Because three of these (Cr, Mn, Ni) are freed in real XRISM cluster
analyses while the other four are weak and tied to iron, we test directly whether
the floor biases the published quantities, and above what exposure.

\paragraph{Design.} We generate one realistic Perseus spectrum from the
independent \code{SpexTruthModel} (all 28 currently available elements at a
literature abundance pattern) and fit it back with the emulator under the
\emph{literature modelling strategy} of the XRISM/Resolve Perseus-core enrichment
analysis \citep{xrism_perseus}: free Si, S, Ar, Ca, Cr, Mn, Fe and Ni together
with $kT$ (or a Gaussian DEM), $\sigma_v$, $\nh$ and the normalisation, with every
other metal tied to the \emph{fitted} iron (a paramless
\code{AbundanceModel.tie\_const}). Injected values follow that analysis
($\mathrm{Fe}=0.71\,\zsun$, $[X/\mathrm{Fe}]$ near unity, $z=0.017284$,
$\nh=1.36\times10^{21}\,\mathrm{cm^{-2}}$, $\sigma_v=180\kms$, $kT=3.9\kev$), over
the $1.9$--$12\kev$ band with the Fe\,XXV complex ($6.567$--$6.620\kev$) masked and
the XRISM/Resolve response ($\sim0.5\,\mathrm{eV}$ channels). Truth and emulator
span the \emph{identical} element set, so their only difference is emulator error.
Code: \code{inference\_demo/hot\_floor/} (\code{experiment.py} builds the store and
a memory-frugal per-element streaming truth; \code{fisher\_bias.py} the estimator
below; \code{plot\_bias.py}, \code{mcmc\_check.py} the figure and cross-check).

\paragraph{Bias versus statistical error.} Since the emulator error is small, the
maximum-likelihood offset in the infinite-count limit is one Newton step from the
truth $\theta_0$,
\begin{equation}
b_{\rm sys}=F^{-1}s,\quad
s_i=\sum_c\Big(\frac{d_c}{\mu_c}-1\Big)\frac{\partial\mu_c}{\partial\theta_i},\quad
F_{ij}=\sum_c\frac{1}{\mu_c}
       \frac{\partial\mu_c}{\partial\theta_i}\frac{\partial\mu_c}{\partial\theta_j},
\end{equation}
with $d$ the noise-free truth counts and $\mu=\mu(\theta_0)$ the emulator, both
restricted to the fit band; $F$ is the Poisson Fisher matrix, so
$\sigma_{\rm stat}=\mathrm{diag}(F^{-1})^{1/2}$ at the reference counts and
$\propto N^{-1/2}$ thereafter. For shape parameters $b_{\rm sys}$ is independent
of the overall count scale (the normalisation absorbs it), and the two quantities
share a single finite-difference Jacobian ($2n{+}1$ forward evaluations). The
crossover
\begin{equation}
N_\star = N_{\rm ref}\,\big(\sigma_{\rm stat}/|b_{\rm sys}|\big)^2
\end{equation}
is the in-band count level at which the emulator bias equals the Poisson error:
below $N_\star$ the floor is invisible, above it it dominates.

\paragraph{Residual character.} The unweighted in-band residual $d/\mu-1$ has
RMS $\sim14\%$, but this lives entirely in the faint inter-line channels (tiny
$\mu$); the \emph{counts-weighted} RMS, which drives the likelihood, is
$0.077\%$---below the $0.1\%$ target---and the noise-free Poisson deviance per
degree of freedom at $10^5$ in-band counts is $3\times10^{-6}$, i.e.\
statistically invisible. This counts-weighted residual is nearly binning
independent (it is the coherent amplitude floor, not a sub-resolution centroid
artifact), confirming the small-residual premise of the linearisation.

\paragraph{Results.} Table~\ref{tab:hotfloor} lists $N_\star$ per parameter for
both temperature models; the fractional systematic biases are $\lesssim0.5\%$
throughout (single-$T$: Fe and the science elements Cr/Mn/Ni at
$0.18$--$0.92\%$). A Perseus \emph{region} carries $\sim4\times10^4$ in-band
counts and a deep stacked core $\sim10^6$; every $N_\star$ exceeds these by
factors of $\sim2\times10^1$ to $\sim10^5$. Equivalently, at a realistic region
the largest emulator-induced bias (Fe and $\nh$) is $\sim4$--$5\%$ of the
$1\sigma$ statistical error, and the hot science elements Cr, Mn and Ni are
$0.1$--$1\%$ of it. The mechanism is exactly the multiplicative character of the
floor: the abundance and normalisation parameters absorb it, and the
worst-emulated metals (Ti, Co, Cu, Zn), being weak and tied to iron, leak
essentially nothing into the free parameters. The Gaussian-DEM fit gives the same
picture (smallest $N_\star=$ Fe at $1.9\times10^7$), so multi-temperature mixing
does not expose the floor. We therefore find that the $0.1\%$ floor miss does
\emph{not} bias XRISM Perseus science at any realistic exposure.

\begin{table}[t]
\centering
\caption{Crossover counts $N_\star$ for the literature-strategy Perseus fit, for
both a single-temperature and a Gaussian-DEM model. $N_\star$ is the in-band
count level at which the emulator systematic bias equals the $1\sigma$ Poisson
error; fractional biases $|b_{\rm sys}|/\mathrm{truth}$ are $\lesssim0.5\%$ for
every parameter in both cases. Realistic anchors: $\sim4\times10^4$ counts per
XRISM region, $\sim10^6$ for a deep stacked core -- every $N_\star$ is orders of
magnitude larger.}
\label{tab:hotfloor}
\begin{tabular}{llcc}
\toprule
Parameter & Group & $N_\star$ single-$T$ & $N_\star$ DEM \\
\midrule
Cr & science-critical & $3.6\times10^{9}$ & $8.6\times10^{9}$ \\
Mn & science-critical & $6.2\times10^{8}$ & $2.6\times10^{9}$ \\
Ni & science-critical & $1.5\times10^{8}$ & $1.1\times10^{8}$ \\
Fe & iron            & $2.6\times10^{7}$ & $1.9\times10^{7}$ \\
Si, S, Ar, Ca & $\alpha$ & $1.2$--$6.3\times10^{8}$ & $0.6$--$17\times10^{8}$ \\
$kT$ / $T_{\rm mean}$ & thermal & $3.5\times10^{9}$ & $2.4\times10^{8}$ \\
$T_\sigma$ (DEM) & thermal & --- & $3.4\times10^{7}$ \\
$\sigma_v$ & nuisance & $4.5\times10^{7}$ & $3.0\times10^{7}$ \\
$\nh$     & nuisance & $1.8\times10^{7}$ & $1.1\times10^{8}$ \\
$Y$ (norm) & nuisance & $3.6\times10^{7}$ & $4.3\times10^{7}$ \\
\bottomrule
\end{tabular}
\end{table}

\begin{figure}[t]
\centering
\includegraphics[width=0.72\textwidth]{../inference_demo/hot_floor/results/bias_vs_counts.png}
\caption{Systematic emulator bias relative to the Poisson statistical error,
$|b_{\rm sys}|/\sigma_{\rm stat}$, versus in-band counts for each fitted
parameter (left: single-$T$; right: Gaussian DEM). Each line has slope $1/2$ and crosses unity
(``bias $=$ noise'', dotted) at $N_\star$; vertical guides mark a realistic XRISM
region and a deep stacked core. All parameters sit far below unity at realistic
exposures.}
\label{fig:hotfloor}
\end{figure}

\paragraph{Caveats and status.} The result is broadening dependent: the
counts-weighted residual rises from $0.077\%$ at $\sigma_v=180\kms$ to $5.4\%$ at
$\sigma_v=0$, so cold, quiet sources are the harder case and are being mapped with
a dedicated low-broadening stress run. The biases above are the linearised
(infinite-count) limit; a direct \code{emcee} counts scan near $N_\star$---where
the systematic offset reaches $\sim1\sigma$ and becomes directly detectable---is
carried out on the cluster (\code{mcmc\_check.py}; the Poisson likelihood cost is
independent of the simulated counts) and is reported next. Both single-$T$ and
Gaussian-DEM Fisher results are reported above and agree; the DEM emulator
temperature grid is held within the training range ($\geq0.7\kev$) to avoid
PCHIP extrapolation of the SPEX truth.

\paragraph{MCMC cross-check.} The linearised bias was validated against a direct
\code{emcee} fit at three in-band count levels ($4\times10^4$, $10^6$, $10^8$),
with a walker-batched GPU forward (\code{gpu\_forward.py}, \code{emcee}
\code{vectorized}). The run uses the production acceleration (TF32 tensor-core
matmul $+$ float32 continuum FFT, ${\sim}3\times$ faster with a ${\sim}10^{-3}$
accuracy cost, below the emulator's own error). All chains are stationary and
unimodal ($\tau\!\approx\!35$--$44$, ESS ${\sim}800$), and the
posterior-predictive residuals are consistent with $\mathcal N(0,1)$ with no
structure at the line complexes. Fig.~\ref{fig:crosscheck} overlays the recovered
offsets $(\mathrm{median}-\mathrm{truth})/\sigma$ on the analytic prediction
$|b_{\rm sys}|/\sigma_{\rm stat}(N)$. At the realistic count levels the predicted
bias is $0.01$--$0.25\sigma$ while the observed offsets are pure statistical
scatter (${\sim}1\sigma$, coverage-consistent): the emulator floor is invisible,
exactly as $N_\star$ predicts. At $10^8$ counts---orders of magnitude beyond any
real observation---the bias emerges, in the \emph{predicted direction}: the sign
of the recovered offset matches $b_{\rm sys}$ for $10$ of the $11$ parameters. Its
\emph{magnitude} runs a factor ${\sim}1.5$--$3$ above the first-order
$b_{\rm sys}$. A dedicated float64 run (TF32/FFT acceleration off) reproduces the
accelerated offsets to ${\lesssim}0.3\sigma$ (Fig.~\ref{fig:crosscheck}, open
squares), so the excess is \emph{not} an acceleration artefact---it is the
first-order linearisation underestimating the true maximum-likelihood bias
(compounded by single-realisation scatter). The implied true $N_\star$ are then a
factor of a few below Table~\ref{tab:hotfloor}, but still ${\gtrsim}10^7$ in-band
counts, and the science-critical Cr/Mn/Ni retain small offsets even at $10^8$
(Cr $-0.2\sigma$, Mn/Ni ${\sim}-1.5\sigma$): the conclusion is unchanged, and the
production acceleration is confirmed science-safe (${\lesssim}0.3\sigma$ at
$10^8$). A multi-seed scan would pin the linearisation factor precisely.

\begin{figure}[t]
\centering
\includegraphics[width=0.68\textwidth]{../inference_demo/hot_floor/results/crosscheck_overlay.png}
\caption{MCMC cross-check: recovered posterior offsets
$|\mathrm{median}-\mathrm{truth}|/\sigma$ (points) at three in-band count levels,
overlaid on the analytic linearised-bias prediction (lines) from
Table~\ref{tab:hotfloor}. At realistic counts (vertical guides) the offsets are
statistical scatter far above the negligible predicted bias; at $10^8$ counts the
bias emerges in the predicted direction (with magnitude a factor of a few above
the first-order lines). Open squares are the float64 reference run at $10^8$; they
coincide with the accelerated points, confirming the acceleration is negligible.
The science-critical Cr/Mn/Ni and Fe are emphasised.}
\label{fig:crosscheck}
\end{figure}

\paragraph{Outlook: NewAthena.} The $0.1\%$ floor is comfortably safe for XRISM,
whose bright-cluster spectra ($\sim4\times10^4$ counts per region, $\sim10^6$ in a
deep stack) lie one to two orders of magnitude below every $N_\star$. NewAthena's
X-IFU, however, collects ${\sim}6\times$ (at $7\kev$) to ${\sim}45\times$ (at
$1\kev$) the XRISM/Resolve effective area \citep{barret2025}, i.e.\ ${\sim}10\times$
the in-band counts for the same source and exposure, so a deep observation of a
bright cluster core reaches ${\sim}10^6$--$10^7$ in-band counts
\citep{cucchetti2018}. That is precisely the regime where the floor first bites:
at ${\sim}10^7$ counts the continuum- and normalisation-sensitive parameters (Fe,
$\nh$, $kT$, $Y$) would carry a ${\sim}0.5$--$1\sigma$ emulator bias, while the
hot-element abundances themselves (Cr/Mn/Ni) require ${\sim}10^8$ and are affected
only in the most extreme stacked spectra. X-IFU's sharper resolution
(${\sim}3$--$4\,\mathrm{eV}$ vs.\ XRISM's ${\sim}4.5\,\mathrm{eV}$) marginally
increases line-head exposure, though the dominant broadening is the cluster's own
${\sim}180\kms$ turbulence, common to both instruments. The floor is therefore a
\emph{trackable systematic} for NewAthena's flagship high-count cluster science---
motivating a line-head retraining of the hot Fe-peak elements before flight---even
though it is negligible for XRISM. These are order-of-magnitude scalings from the
effective-area ratio; a dedicated X-IFU/\code{SIXTE} simulation would pin $N_\star$
in X-IFU counts.

\section{Bias across parameter space}\label{sec:biassweep}
The study of Sec.~\ref{sec:hotfloor} evaluates $b_{\rm sys}$ and $N_\star$ at a
\emph{single} point in parameter space---Perseus---plus a small number of
single-axis variants. That is the right design for the question it asks (does
the hot-element floor damage XRISM Perseus science?), but it does not establish
that the emulator is unbiased for cluster science generally. The estimator of
Sec.~\ref{sec:fisherbias} is cheap enough to answer the general question
directly, by evaluating it over a sample of the parameter space rather than at
one point.

\paragraph{Design.} We draw a Latin hypercube over the cluster science range:
$kT$ (or $T_{\rm mean}$) $1.5$--$8\kev$, $T_{\rm sigma}$ $0.15$--$3.0\kev$ for
the DEM flavour, each of the eight free abundances independently over
$0.2$--$2\,\zsun$, $\sigma_v$ $30$--$600\kms$, and $\nh$ $0$--$5\times10^{21}\,
\mathrm{cm^{-2}}$. A Latin hypercube rather than a grid: with 12--13 axes a grid
is combinatorially impossible, while the hypercube gives every axis full
marginal coverage at any sample size, so a partially completed run is still
interpretable. The range is deliberately the volume real cluster fits occupy
rather than the full training box---filling the map with temperatures no cluster
exhibits would swamp the science-relevant signal.

\paragraph{Why the abundance axes matter.} Abundance enters both the truth and
the emulator \emph{linearly} (CIE emissivities sum per element), so varying it
cannot by itself create per-element emulator error. It nonetheless moves
$b_{\rm sys}$, because it changes which elements dominate which channels and
therefore reweights the existing per-element errors in the score,
Eq.~\eqref{eq:score}. This is an effect that per-element accuracy metrics
cannot express: it lives in the parameter covariance $F^{-1}$, not in any
single element's residual. It is the main reason the sweep varies abundances
rather than holding them at the fiducial pattern.

\paragraph{Implementation.} \code{scripts/bias\_sweep.py}, in two resumable
stages matching the split of Sec.~\ref{sec:fisherbias}: \code{--stage truth}
computes $d$ at every sample point, and \code{--stage bias} computes $\mu$, $J$
and hence $b_{\rm sys}$ and $\sigma_{\rm stat}$. The stages are separate because
$d$ is a property of SPEX alone: when an element is retrained or the store is
rebuilt, only the second stage is repeated. The truth stage loops
\emph{elements outermost}, loading each element's cache once and evaluating it
at every sample point, which keeps peak memory at a single element
(${\sim}0.4$\,GB) and avoids re-reading the caches once per point.

\paragraph{Interpretation.} Results are reported per parameter as the
distribution of $|b_{\rm sys}|/\sigma_{\rm stat}(N)$ over the sampled volume at
a chosen exposure $N$, together with the median $N_\star$ and the fraction of
sampled points exceeding $1\sigma$. Following Sec.~\ref{sec:fisherbias} these
are a screen: points flagged above $1\sigma$ are candidates for confirmation
with full posterior fits, where the true offset may be larger by a factor of a
few.

\section{Sampling the velocity dispersion}\label{sec:sigmav}
The line-of-sight velocity dispersion $\sigma_v$ is a free, sampled parameter.
It was previously pinned in the walker-batched runs for an implementation
reason rather than a physical one: the analytic line deposit
(\code{deposit\_gaussian\_lines}) accepted only a scalar velocity, while the
continuum FFT already broadened per spectrum. The two branches differ in
bookkeeping, not in physics. With one $\sigma_v$ every walker shares a line's
$\pm n\sigma$ window, so the scatter indices are common to the batch; with a
per-walker $\sigma_v$ those windows become ragged across the batch, and the
deposit is instead accumulated into a flattened buffer with walker-offset
indices. The \code{erf}-difference weights, and hence each line's exactness on
any target binning, are unchanged.

The prior is uniform on
\begin{equation}
\sigma_v \in [30,\,600]\kms ,
\end{equation}
grounded in the measured Perseus values and then widened. \emph{Hitomi} found
$\sigma_v = 164\pm10\kms$ at $30$--$60\,$kpc from the nucleus
\citep{hitomi2016}, and mapping the core gave ${\sim}100\kms$ over most of the
field, rising to ${\sim}200\kms$ toward the central AGN and the north-western
ghost bubble, with a ${\sim}100\kms$ line-of-sight gradient from sloshing
\citep{hitomi2018}. \emph{XRISM}/Resolve subsequently measured ${\sim}300\kms$
in the eastern region and a $\pm200$--$300\kms$ dipole attributed to a recent
merger, out to ${\sim}500\,$kpc \citep{xrism_perseus_kinematics}. Perseus itself
therefore spans ${\sim}100$--$300\kms$; the prior is widened on both sides so
the boundary is never informative. The floor is kept strictly positive because
as $\sigma_v\to0$ the line profile collapses below the instrumental response
width and the parameter ceases to be identifiable --- $30\kms$ is already well
below Resolve's own resolution, which corresponds to ${\sim}100\kms$ at Fe-K.
The bias study deliberately retains a wider $[10,1000]\kms$ ``extreme'' range so
that calibration is also probed outside anything the sample actually shows.

\section{Cost of the forward step}\label{sec:forwardcost}
One vectorised MCMC step is one batched forward over the walker ensemble, so its
cost sets the feasibility of any large sampling or simulation-based-calibration
campaign. Measured on an A10-class GPU with the full 30-element store at 96
walkers, evaluating the per-element trunks together (grouped \code{vmap} over
same-shape elements) and then compiling that batched call gives
\begin{center}
\begin{tabular}{lrrrr}
\toprule
configuration & ms/walker & vs.\ fp64 & vs.\ prod. & max.\ rel.\ err. \\
\midrule
serial, float64, no acceleration        & 411.7 & $1.00\times$ & $0.60\times$ & (reference) \\
serial + TF32 + compile + float32 FFT   & 246.2 & $1.67\times$ & $1.00\times$ & $9.7\times10^{-4}$ \\
batched (grouped \code{vmap})           & 208.6 & $1.97\times$ & $1.18\times$ & $6.2\times10^{-4}$ \\
\quad + \code{torch.compile}            & 112.7 & $3.65\times$ & $2.19\times$ & $6.2\times10^{-4}$ \\
\bottomrule
\end{tabular}
\end{center}
The second row is the current production configuration, so the batched and
compiled forward is $2.19\times$ faster than what is deployed, and
$3.65\times$ faster than the unaccelerated reference. Grouping the element
trunks is worth only $1.18\times$ on its own; compiling the resulting grouped
GEMMs is what pays. That compile step is a genuine gain rather than a duplicate
of the per-element compilation applied elsewhere: the batched wrapper
deliberately calls the unbound \code{forward\_norm} so as to bypass any compiled
instance attribute, so before this the batched trunk was never compiled at all.

Accuracy is not traded for any of this. Against the float64 reference, on bins
carrying at least $0.1\%$ of the peak flux, the deployed accelerations cost
$9.7\times10^{-4}$ at worst and the batched path $6.2\times10^{-4}$, with $99$th
percentiles of $2.8$ and $2.4\times10^{-4}$; the compile step changes the result
by less than one part in $10^{3}$ of the error itself. All of this sits well
below the emulator's own ${\sim}3\times10^{-3}$ test error.

The forward is strongly dominated by the trunk. A stage breakdown at batch $16$
gives trunk $95\%$, continuum broadening and rebinning $2\%$, and line
deposition plus abundance weighting $2\%$; an operator-level profile attributes
$63.6\%$ of GPU time to \code{bmm} (executing as the TF32 tensor-core
\code{cutlass\_80\_tensorop\_s1688gemm} kernel) and ${\sim}18\%$ to the
bandwidth-bound activation kernels. The trunk is therefore GEMM-bound and
already on tensor cores.

That leaves little room, and we record two optimisations that were implemented,
measured and \emph{rejected}, because both failures are informative.
\begin{itemize}
\item \textbf{bfloat16 autocast on the trunk} is fast --- $1.69\times$, to
$66.9\,$ms/walker --- but costs a maximum relative error of $2.4\%$ on a bin at
$97.5\%$ of the spectral peak, with a $99$th percentile of $8.7\times10^{-3}$,
roughly three times the emulator's own ${\sim}3\times10^{-3}$ test error. A
systematic of that size in the forward model is not acceptable in a study whose
purpose is to calibrate the posterior.
\item \textbf{Coarsening the trunk energy grid} fails because its premise is
false. The trunk is not a smooth continuum: measured on the native
$P=24212$ grid, the one-bin curvature $|d^2\log_{10}\rho|$ peaks at $0.0016$ for
He but at $1.75$, $2.45$ and $4.91$ for O, Si and Fe respectively. Only H and He
are genuinely smooth; every metal carries line structure in the trunk, including
elements that also have a line head. Evaluating every second point and
interpolating back costs Fe an order of magnitude in density error and $43.5\%$
in broadened flux, with the worst bin at $70\%$ of the peak --- and the
broadening does not wash this out, because narrow lines are being lost rather
than noise being added. The same objection applies to any scheme that
interpolates the trunk over a precomputed temperature grid.
\end{itemize}
Increasing the walker sub-batch changes nothing (the GPU is already saturated at
$16$), and CUDA graph capture is not indicated: the profile shows the CPU
blocked on a full command buffer, i.e.\ running ahead of the GPU rather than
failing to feed it. We therefore treat the forward as converged at
$1.88\times$, and note that at ${\sim}5.9\,$h per chain the productive direction
for a calibration campaign is fewer forward evaluations, not faster ones.

\section{Sampler comparison}\label{sec:bakeoff}
Because the forward dominates the cost of every algorithm
(Sec.~\ref{sec:forwardcost}), the choice of sampler decides whether a
calibration campaign is affordable at all. We therefore fit the \emph{same}
simulated Perseus spectrum ($10^{6}$ in-band counts, 12 free parameters, the
literature modelling strategy of Sec.~\ref{sec:hotfloor}) with each sampler over
an identical \code{PoissonPosterior}, so that what is compared is the algorithm
and nothing else. Code: \code{scripts/bake\_off.py}.

\paragraph{Results.} Table~\ref{tab:bakeoff} gives the state of the comparison.
\emph{These are partial. The flow-accelerated samplers (\code{nautilus},
\code{pocoMC}, \code{i-nessai}) are implemented and verified against analytic
problems but not yet run at scale; NUTS was diagnosed rather than completed;
and the variational row is being repeated with a defensible particle count.
The three ensemble rows are final.}

\begin{table}[t]
\centering
\caption{Sampler comparison on one $10^{6}$-count Perseus spectrum, 12 free
parameters, on an A10-class GPU. ``Evals'' counts walkers pushed through the
emulator. Agreement is against the \code{emcee} chain.}
\label{tab:bakeoff}
\begin{tabular}{lrrrrrl}
\toprule
sampler & wall (h) & evals & min ESS & ESS/eval & ms/eval & agreement \\
\midrule
\code{emcee}     &  1.65 &  51\,020 &  144 & $2.8\times10^{-3}$ & 117 & (reference) \\
\code{zeus}      &  5.38 &  70\,800 &  194 & $2.7\times10^{-3}$ & 274 & $0.21\sigma$, width $1.04$--$1.24$ \\
\code{UltraNest} & 10.83 & 218\,711 & 3507 & $1.6\times10^{-2}$ & 178 & $0.14\sigma$, width $0.94$--$1.11$ \\
\code{SVI}       &  1.25 &   8\,006 & 4000\rlap{$^\dagger$} & --\rlap{$^\dagger$} & 564 & $2.57\sigma$, width $0.43$--$2.96$ \\
\bottomrule
\end{tabular}
\end{table}

\paragraph{ESS per evaluation is not the whole story.} Two rows of
Table~\ref{tab:bakeoff} must not be read as they stand. $^\dagger$A variational
posterior is independent by construction, so its ``ESS'' is simply the number of
draws requested and its ESS/eval is an artefact of two settings rather than a
measurement; the variational row is judged on recovery and agreement only.
Separately, the importance-based samplers report a Kish effective size over
their weights, which \emph{is} comparable as an ESS but corresponds to a much
larger raw sample array. More fundamentally, ESS/eval and wall-clock measure
different things here: the ensemble samplers amortise a whole batch of walkers
into each forward, so their per-evaluation cost is the batched
$112$--$131\,$ms/walker of Sec.~\ref{sec:forwardcost}, whereas a single-point
sampler pays ${\sim}4\times$ that for a $B{=}1$ forward. A method can therefore
need far fewer evaluations and still lose badly on wall-clock; that is a fact
about batching, not about mixing.

\paragraph{UltraNest is the most efficient, but not the cheapest.} Nested
sampling delivers $5.7\times$ more effective samples per evaluation than
\code{emcee} and $3.7\times$ more per wall-clock second, and it is the only
sampler here that returns $\log Z$. Its cost is nonetheless irreducible: it runs
to an evidence tolerance rather than a step count, so it cannot be asked for
fewer draws. It produced ESS $3507$ where simulation-based calibration needs
only ${\sim}100$, at $10.8\,$h. For a campaign of many fits \code{emcee}
remains the economical choice ($1.65\,$h per posterior) with UltraNest as the
reference; the information gain from prior to posterior here is ${\sim}44$ nats,
which sets the nested-sampling iteration count.

\paragraph{Slice sampling is defeated by the hardware, not the algorithm.}
\code{zeus} needs only ${\sim}5.5$ evaluations per walker per iteration -- its
mixing is as advertised -- but its vectorised implementation hands the sampler a
\emph{shrinking} active subset as walkers accept, so the tail of every iteration
runs at $B{=}8,4,2$ where the GPU is idle. Its effective cost is
$274\,$ms/eval against \code{emcee}'s $117$. Combined with posteriors $4$--$24\%$
wider, it is dominated on every axis.

\paragraph{Gradient-based sampling is defeated by the batch size, not by
mixing.} Pyro's NUTS ran at $958\,$s per iteration, saturating its tree-depth
ceiling at $1023$ leapfrog steps, each a gradient at $B{=}1$ costing $937\,$ms
-- implying a single-point forward of ${\sim}526\,$ms, $4\times$ the batched
$131\,$ms/walker. The comparison that matters is not wall-clock but the
break-even: a NUTS gradient costs $7.1\times$ an amortised ensemble evaluation,
so NUTS must reach $\mathrm{ESS/eval}>2.0\times10^{-2}$ to win, and since a
well-adapted NUTS yields roughly one independent draw per iteration its
ESS/eval is ${\sim}1/(\text{leapfrog steps})$. \emph{Break-even is therefore
around $50$ steps}, and at $31$ steps NUTS would beat every other sampler here.
Running at $1023$ it is a factor ${\sim}21$ from viability, which is a statement
about adaptation, not about the algorithm. Two defaults are implicated: Pyro
adapts a \emph{diagonal} mass matrix, which cannot represent correlation at all
and so cannot follow this posterior's abundance/normalisation degeneracy (the
marginal scales are not the difficulty -- their condition number in the
unconstrained space is only ${\sim}79$, worth ${\sim}9$ steps); and the default
initialisation draws from the prior box, whereas every other sampler here starts
at the truth. Both are now options. Should they prove insufficient, the
indicated direction is an \emph{ensemble} gradient method (ChEES-HMC, MEADS)
that evaluates many chains per forward, since the batched potential and its
gradient are already available; at $B{=}64$ a gradient costs $233\,$ms per
chain rather than $936$.

\paragraph{Variational inference is fast and wrong, and the two are related.}
SVI was the quickest ($1.25\,$h) and eleven of its twelve parameters look
healthy (pulls of mean $+0.18$, standard deviation $1.10$, indistinguishable
from \code{emcee}'s $-0.17$, $1.00$). The exception is decisive: $\nh$ is
recovered $8\sigma$ from the truth \emph{in its own units}, which decomposes
into a $2.6\sigma$ offset combined with a posterior $0.43\times$ too narrow. It
is the dangerous kind of failure: confidently wrong error bars on one parameter,
with nothing anomalous elsewhere.

The speed was, however, borrowed from the accuracy. A variational run costs
exactly $(\text{steps})\times(\text{particles})$ forward evaluations, and this
one used only four ELBO particles -- a memory-driven choice, not an
accuracy-driven one -- for $8\,000$ evaluations, a sixth of \code{emcee}'s. Four
samples is a very noisy estimate of the ELBO gradient, and restoring a
defensible particle count costs the speed advantage outright: staying cheaper
than \code{emcee} allows at most ${\sim}25$ particles over $2000$ steps.

We do not yet attribute the failure. The natural suspicion is an inadequate
variational family, but the evidence points the other way: minimising the
reverse Kullback--Leibler divergence is mode-seeking and systematically
\emph{under}-estimates variance, whereas several of these posteriors came out up
to $2.96\times$ \emph{too wide}, which is the signature of an unconverged
optimiser rather than of a family mismatch. A full-rank Gaussian can in any case
represent a linear degeneracy exactly, and at $10^{6}$ counts with ${\sim}44$
nats of information the posterior should be close to Gaussian -- an assumption
the Fisher treatment of Sec.~\ref{sec:fisherbias} already relies on and
validates. The informative next experiment is therefore not a more expressive
guide but a better-initialised one: $F^{-1}$ from Sec.~\ref{sec:fisherbias}
\emph{is} the posterior covariance, so initialising the guide's location at the
truth and its scale at $\mathrm{chol}(F^{-1})$ starts the optimisation with the
correct correlation structure instead of requiring it to be discovered by
stochastic gradient descent. Only if that still fails is a normalizing-flow
guide indicated -- and it carries $2\,786$ parameters against the full-rank
Gaussian's $168$, so it would need proportionally more steps and make the cost
problem worse.

\paragraph{Recovery is consistent across samplers, and consistent with noise.}
All three sampling methods agree on the same parameter offsets --
$\mathrm{Ar}$ at ${\sim}-2\sigma$, $\mathrm{Ca}$ at ${\sim}+1.4\sigma$ -- and
agree with each other to $0.21\sigma$ (zeus) and $0.14\sigma$ (UltraNest). Those
offsets are therefore a property of the simulated data, not of the samplers. They
are also fully consistent with counting noise: across the 12 parameters the
\code{emcee} pulls have mean $-0.23$ and standard deviation $0.98$ against
expectations of $0\pm0.29$ and $1$, giving $\chi^{2}=11.1$ for 12 degrees of
freedom ($p=0.52$). The Fisher estimator of Sec.~\ref{sec:fisherbias} predicts a
maximum emulator-induced bias of $0.23\sigma$ at this exposure, twenty times
smaller than the largest observed pull. A single fit therefore cannot measure
emulator bias, which is precisely the argument for the sweep of
Sec.~\ref{sec:biassweep} and for calibration over many realisations rather than
one showcase.


\begin{thebibliography}{9}
\bibitem[Morrison \& McCammon(1983)]{morrison1983} Morrison, R.\ \& McCammon, D.\
1983, ApJ, 270, 119.
\bibitem[Wilms et al.(2000)]{wilms2000} Wilms, J., Allen, A.\ \& McCray, R.\
2000, ApJ, 542, 914.
\bibitem[XRISM Collaboration(2026)]{xrism_perseus} XRISM Collaboration 2026,
Chemical enrichment of the Perseus cluster core seen by XRISM/Resolve,
arXiv:2606.17141.
\bibitem[Hitomi Collaboration(2016)]{hitomi2016} Hitomi Collaboration 2016,
Nature, 535, 117 (The quiescent intracluster medium in the core of the Perseus
cluster).
\bibitem[Hitomi Collaboration(2018)]{hitomi2018} Hitomi Collaboration 2018,
PASJ, 70, 9 (Atmospheric gas dynamics in the Perseus cluster observed with
Hitomi).
\bibitem[Zhang et al.(2025)]{xrism_perseus_kinematics} Zhang, C., et al.\ 2025,
Mapping the Perseus galaxy cluster with XRISM: gas kinematic features and their
implications for turbulence, arXiv:2510.12782.
\bibitem[Cucchetti et al.(2018)]{cucchetti2018} Cucchetti, E., Pointecouteau, E.,
Peille, P., et al.\ 2018, A\&A, 620, A173 (Athena X-IFU synthetic observations of
galaxy clusters to probe the chemical enrichment of the Universe).
\bibitem[Barret et al.(2025)]{barret2025} Barret, D., et al.\ 2025, The X-ray
Integral Field Unit at the end of the Athena reformulation phase, arXiv:2502.10866.
\end{thebibliography}

\end{document}
